\documentclass[11pt]{article}
\usepackage[margin=1in]{geometry}
\usepackage{amsmath}
\usepackage{booktabs}
\usepackage{graphicx}
\usepackage{hyperref}
\usepackage{natbib}
\usepackage{todonotes}

\title{Does Relational Information Improve Stock Return Ranking?\\
\large A Replication of Temporal Relational Ranking for Stock Prediction}
\author{}
\date{}

\begin{document}
\maketitle

\begin{abstract}
\todo[inline]{Write last, once Sections 5--7 have final numbers.}
\end{abstract}

\section{Introduction}
\label{sec:intro}

A stock's future return is usually modeled as a function of its own history: its own past
prices, its own past volumes, its own fundamentals. This is convenient, but it discards
information that is easy to observe and often economically meaningful --- companies in the same
industry, connected by a supply chain, or held together in the same funds tend to move together,
because a shock to one is frequently a shock to the group. A model that treats every stock as an
independent time series is, by construction, blind to that channel.

\citet{feng2019temporal} address this directly. Their Relational Stock Ranking (RSR) model
encodes a graph of relations between stocks --- who shares an industry, who is connected through
Wikidata facts about ownership, supply, and competition --- into a Temporal Graph Convolution
layer, and combines it with a sequential (LSTM-based) encoding of each stock's own price history.
The relational signal is used to re-weight and propagate information between connected stocks at
every training step, and the model is trained to rank stocks by next-day return rather than to
regress the return value directly, since for a long-short or top-$k$ trading strategy the relative
ranking of stocks matters more than the exact predicted return.

This project replicates RSR on the same NASDAQ universe and industry-relation graph used in the
original paper, asking whether its central empirical claim --- that encoding relational
information improves next-day return ranking relative to a relation-blind baseline --- holds up
when the model is retrained from scratch on the authors' own released data and code.

Our interest in this specific question comes from outside this course: we separately maintain a
multi-agent system for trading and investment decisions, in which individual agents each analyze
a single stock along a different dimension (technical signals, news/sentiment, risk) before their
outputs are aggregated into a decision. That system currently has no agent whose input is the
relationship between stocks --- every agent's view of a stock is, in effect, independent of every
other stock, the same simplifying assumption \citet{feng2019temporal} argue against. RSR is a
concrete, reproducible instance of turning ``stocks are related'' into a model input (a relation
graph plus a learned attention weight over neighbors) rather than an informal claim, which makes
it a useful first test of whether that kind of signal is worth adding to such a system. This
project does not attempt that integration or extend RSR's architecture; it replicates RSR on its
own terms, as a first check of whether the relational signal it relies on holds up before any
downstream use is considered.

\subsection{Research question}

\begin{itemize}
  \item \textbf{RQ1 (replication).} Does encoding a relation graph between stocks, following
  \citet{feng2019temporal}, improve next-day return ranking relative to a baseline that treats
  stocks independently?
\end{itemize}

\subsection{Contribution}

We replicate the RSR baseline and relational model on the official NASDAQ dataset, training both
to convergence under a matched schedule. RQ1 does not hold in the direction the original paper
reports: at convergence, encoding the static industry relation graph does not lower mean-squared
prediction error relative to the no-relation baseline, and the two models split on the two
ranking-adjacent metrics --- the baseline has the higher top-1 ranking accuracy (mrrt), while the
relational model has the higher simulated trading return (btl). We report this split result, and
the environment and data issues we had to resolve to obtain it (Sections~\ref{sec:data} and
\ref{sec:methods}), as the project's substantive findings.

\section{Related Work}
\label{sec:related}

\textbf{Relation-based stock prediction.} \citet{feng2019temporal} is the paper this project
replicates; Section~\ref{sec:methods} describes its architecture in the detail needed to reproduce
it. \citet{cui2023temporal} extend the same broad idea --- that a stock's neighbors carry
predictive information --- from pairwise relations to group-wise hypergraph relations
(industry-belonging and fund-holding hyperedges), with a hierarchical attention mechanism over
intra-hyperedge, inter-hyperedge, and inter-hypergraph structure. \citet{jafari2022gcnet} instead
use a graph convolutional network directly on a stock-relation graph for price-movement
classification; the authors note in their code release that the implementation was withheld
pending journal review, so we treat this as a design reference rather than a reproducible
baseline. All three approaches share the same core assumption this project tests directly: that
encoding a graph of relations between stocks improves prediction relative to treating each stock
independently.

\section{Data}
\label{sec:data}

\subsection{Source}

We use the NASDAQ portion of the dataset released with \citet{feng2019temporal}
(\texttt{Temporal\_Relational\_Stock\_Ranking}, official repository), which contains: (i) daily
end-of-day features for 1{,}026 NASDAQ tickers --- 5-, 10-, 20-, and 30-day moving averages of
close price plus the raw close-to-close return ratio, already normalized by the stock's maximum
price over the sample --- and (ii) a static industry-relation graph over the same 1{,}026 tickers,
a one-hot encoding of shared-industry membership across 97 industry categories sourced from
sector/industry classifications.

The repository's own documentation describes the underlying raw pull as ``historical (30 years)
End-of-day data \ldots{} of more than 8,000 stocks.'' That describes the unfiltered Google Finance
pull, not the modeling dataset actually used for training: the processed feature files cover
\textbf{1{,}246 trading days per stock, roughly five years}, not thirty, and only 1{,}026 of the
raw universe's stocks survive the paper's own liquidity/history filtering. A reader relying on the
top-level description alone would substantially overestimate both the historical depth and the
breadth of the modeling sample.

\subsection{Data quality}

Exploratory analysis (full notebook: \texttt{reports/R2\_data\_and\_eda/eda.ipynb}) surfaced two further
gaps between the data and its documentation, beyond the sample-size point above:

\begin{enumerate}
  \item \textbf{Missing values are not \texttt{NaN}.} The preprocessing pipeline
  (\texttt{preprocess/eod.py}) pads missing trading days with a sentinel value of $-1234$ rather
  than a null. Undetected, this sentinel dominates naive summary statistics: raw skewness for
  every feature column comes out around $-23$, which would be an extraordinary departure from
  normality for a bounded price ratio. Once the sentinel is stripped, skewness for every column
  falls to an unremarkable $-0.3$ to $-0.4$. The original $-23$ number says nothing about the
  price data; it is an artifact of an undocumented missing-value code.
  \item \textbf{The relation graph includes non-company tickers.} 156 of the 1{,}026 NASDAQ
  tickers (15.2\%) are recorded with industry \texttt{"n/a"} in the sector/industry relation file.
  Inspecting these tickers shows they are ETFs and index funds (e.g.\ \texttt{QQQ}, \texttt{SHV},
  \texttt{IEF}, \texttt{PFF}) rather than operating companies. An ETF's return is a basket average
  with no industry relation of its own, so its presence in a graph meant to encode
  company-to-company relations is noise rather than signal; 17 of these 156 tickers are fully
  isolated (degree zero) in the industry graph as a direct consequence.
\end{enumerate}

Both issues are corrected identically across every model reported in Section~\ref{sec:methods}:
the $-1234$ sentinel is treated as missing (not as a value) wherever it appears, and results are
reported with the 156 ETF/fund tickers still present in the universe unless noted otherwise, since
removing them would change the size of the relation graph and is left as a robustness check in
\texttt{reports/R3\_model\_and\_experiments/PROGRESS.md}.

\section{Methods}
\label{sec:methods}

We train and compare two models, both sharing the same NASDAQ universe, the same 5-, 10-, 20-,
30-day moving-average and close-ratio features, and the same next-day ranking objective. Both
are trained with the same optimizer (Adam), the same train/validation/test split by date
(indices 0--755 / 756--1007 / 1008-- of the 1{,}246 trading days), and the same combined loss:
a mean-squared-error regression term on predicted return plus a pairwise ranking term that
penalizes any pair of stocks whose predicted return order disagrees with their realized order,
weighted by $\alpha$.

\subsection{Baseline: Rank\_LSTM}

The baseline (\texttt{training/rank\_lstm.py} in the official repository) encodes each stock's own
feature history with a single-layer LSTM and predicts next-day return from the final hidden state,
with no information exchanged between stocks. This is the ``stocks are independent'' benchmark
that Section~\ref{sec:intro} argues is the wrong assumption, and the lower bound RSR is designed to
beat.

\subsection{RSR: static relation graph}

The RSR model (\texttt{training/relation\_rank\_lstm.py}) takes the same per-stock representation
--- here, the baseline LSTM's learned sequential embedding, exported after training rather than
downloaded, since the original authors' pretrained embedding file is distributed only via an
external link that we chose not to depend on --- and propagates it across the static industry
relation graph via a learned, softmax-normalized attention weight over each stock's neighbors,
before combining the propagated and original representations into the final prediction. The
relation graph itself, a $1{,}026\times1{,}026\times1$ one-hot encoding, is a training-time
constant: identical on day one of training and day 1{,}246.

\section{Experiments and Results}
\label{sec:results}

Both models are trained with $\text{seq}=4$, $\text{unit}=32$, the same Adam optimizer, and
the same combined MSE-plus-ranking loss described in Section~\ref{sec:methods}, on the same
NASDAQ universe (1{,}026 tickers, Section~\ref{sec:data}), for the same 50-epoch schedule. We
report three metrics on the held-out test split: \emph{MSE}, the mean-squared error of predicted
return; \emph{mrrt}, the mean reciprocal rank of the true top-return stock (higher means the
model's top pick is more often the stock that actually gained the most); and \emph{btl}, the
simulated return of a strategy that buys the model's top-ranked stock each day (a simple
long-only backtest).

\begin{table}[h]
\centering
\caption{Test-set performance, both models trained to 50 epochs on the same NASDAQ universe
and embedding.}
\label{tab:results50}
\begin{tabular}{lccc}
\toprule
Model & Test MSE & Test mrrt & Test btl \\
\midrule
Baseline (Rank\_LSTM, no relation) & 0.0003774 & \textbf{0.0488} & 1.018 \\
RSR (static industry relation) & 0.0003774 & 0.0276 & \textbf{1.130} \\
\bottomrule
\end{tabular}
\end{table}

Table~\ref{tab:results50} shows two results that both cut against the original paper's headline
claim.

First, \textbf{MSE is essentially identical across the two models} --- the baseline and RSR agree
to seven decimal places. Encoding the static industry relation graph does not lower mean-squared
prediction error relative to treating stocks independently once training is run to convergence:
RQ1, in the specific sense of ``does a relation graph reduce MSE,'' is not supported here.

Second, \textbf{the two ranking-adjacent metrics point in opposite directions}. The no-relation
baseline has the higher mrrt (0.0488 vs.\ 0.0276), meaning it more often puts the single best-
performing stock in first place; RSR has the higher simulated trading return (1.130 vs.\ 1.018),
meaning a strategy that simply buys its top pick each day does better over the test period. Which
model ``wins'' on RQ1 therefore depends on which of these two metrics one treats as the target: a
strategy judged purely by whether it nails the single top pick favors the baseline, while a
strategy judged by realized trading return over the test period favors RSR.

\section{Discussion}
\label{sec:discussion}

\subsection{Did the replication succeed?}

Not in the direction \citet{feng2019temporal} report. The original paper's central claim is that
encoding relational information lowers prediction error and improves ranking relative to a
relation-blind baseline; at convergence, our replication finds neither (Table~\ref{tab:results50},
baseline vs.\ RSR rows). We see three non-exclusive reasons this replication may diverge from the
original result, none of which we can rule in or out with a single training run:

\begin{enumerate}
  \item \textbf{Hyperparameters were not tuned to this setting.} We used a single configuration
  ($\text{seq}=4$, $\text{unit}=32$) chosen for consistency between the baseline embedding and the
  relational model, not a value selected by the kind of hyperparameter search the original paper
  likely performed. The original paper's own example commands use larger sequence lengths and
  hidden sizes (e.g.\ $\text{seq}=16$, $\text{unit}=64$ for their published wikidata run).
  \item \textbf{A single training run.} Every number in Table~\ref{tab:results50} comes from one
  run per model with a fixed random seed. A single 50-epoch run is not by itself strong evidence
  that training has fully converged or that the ranking between models is stable; it is one draw,
  not a distribution.
  \item \textbf{We used the pretrained embedding, but exported it ourselves.} Because the
  original authors' pretrained sequential embedding is only available via an external link
  (Section~\ref{sec:methods}), the embedding RSR consumes here was exported from our own baseline
  run rather than the authors' released artifact, which may not be trained under the same
  conditions as the one used in the original paper's reported numbers.
\end{enumerate}

We do not have evidence to prefer one of these explanations over the others; distinguishing them
is exactly what the robustness checks in \texttt{reports/R3\_model\_and\_experiments/PROGRESS.md} (a second
random seed; a brief hyperparameter sweep) are for, and that work is not yet done as of this
draft.

\subsection{Limitations}

Results are reported for NASDAQ only, using the industry relation graph only (the original paper
also reports a Wikidata-relation condition, which we did not run); the 156 ETF/fund tickers
identified in Section~\ref{sec:data} remain in the relation graph for both models reported here,
rather than being excluded as their absence of a true industry relation would suggest; and only
one ablation dimension (relation graph present vs.\ absent) was tested, on one market.

\section{Conclusion and Future Work}
\label{sec:conclusion}

This project replicated \citet{feng2019temporal}'s Relational Stock Ranking model on its official
NASDAQ dataset, comparing it against a no-relation Rank\_LSTM baseline under a matched training
schedule (50 epochs, same optimizer, same embedding). At convergence, encoding the static industry
relation graph did not lower mean-squared prediction error relative to the relation-blind baseline
--- the two models agree on MSE to seven decimal places --- and the baseline in fact had the higher
top-1 ranking accuracy (mrrt), while RSR had the higher simulated trading return (btl). The
replication therefore does not reproduce the original paper's headline claim that relational
information improves both prediction error and ranking performance in this setting; the direction
depends on which of the two ranking-adjacent metrics (mrrt vs.\ btl) one weights more.

Following the course's framing of ``potentials of your projects,'' the most direct next steps are
the remaining robustness checks scoped in \texttt{reports/R3\_model\_and\_experiments/PROGRESS.md} --- a second
random seed and a brief hyperparameter sweep closer to the original paper's published
configuration --- to establish whether the mrrt/btl split and the absence of an MSE improvement
are stable findings or artifacts of a single run. Extending the comparison to the Wikidata relation
and to NYSE, both available in the same official dataset but not run here, is the natural next
step after that.

\bibliographystyle{chicago}
\bibliography{references}

\end{document}
